\subsection{Resumen de la charla}
\textit{Nombre de la platica}: Apertura a la descentralización.
\vspace{1cm} 

Muchos grandes modelos de lenguaje (LLMs) se presentan como abiertos, pero en realidad no lo son.
La apertura de los LLMs a menudo se limita a la capacidad de acceder al modelo sin la receta de como se construyó, ni sus datos
de entrenamiento; por ello se les denomina \textit{modelos de peso abierto}. 
Actualmente, la mayoría de los LLMs locales son solo modelos de peso abierto, aunque existen algunos, como OLMo\cite{groeneveld-etal-2024-olmo} o Pythia\cite{biderman2023pythia}, 
que son completamente abiertos. Esto dificulta rastrear el origen del modelo, sus datos de entrenamiento -necesario para rastrar denucias de derechos de autor- y sus limitaciones, ocasionando una falta de transparencia que puede generar 
problemas de responsabilidad, sesgo y uso indebido.

En esta charla presentaremos algunas legislaciones y regulaciones que buscan la apertura de los LLMs, y argumentaré que la verdadera apertura de los LLMs implica no solo el acceso a los pesos, sino también a los datos de entrenamiento y a la capacidad de modificar y auditar el modelo. La discusión sobre LLMs “abiertos” no es solo técnica, sino una disputa sobre quién controla los medios de producción cognitivos, quién define la verdad y quién conserva la capacidad de deliberación y autonomía. La charla está organizada en cuatro ejes de discusión: 1) infraestructura (Marx), 2) control (Foucault), 3) conocimiento (Habermas, esfera pública) y 4) relaciones (Serapio: parasentimentalidad).